\section{Results}
\label{sec:results}

\paragraph{Main result.} Table~\ref{tab:test} gives the test-set results. Prompted directly, none of the three local
models can turn the training siblings into a working rule: Qwen2.5-Coder solves \testQwenCoderOneShot{} of
\testBuildable{} conversions and the other two none. With the framework, Qwen2.5-Coder solves
\testQwenCoderFull{} on average (\testQwenCoderFullMin--\testQwenCoderFullMax{} across seeds), Llama~3.1 8B
solves \testLlamaFull{} and Gemma~3 4B \testGemmaFull. Spinfer solves \testSpinfer{} and the frontier LLM agent
\testFrontier. The local model therefore closes most but not all of the distance to Spinfer, and the frontier agent
stays well ahead. The per-migration picture is uneven: the guided Qwen2.5-Coder solves more conversions than
Spinfer on \texttt{pcim\_request\_all\_regions} and \texttt{secs\_to\_jiffies}, matches it on
\texttt{iio\_claim\_direct}, and falls behind on \texttt{iio\_push\_with\_ts}; on \texttt{gpio\_set\_rv} and
\texttt{ethtool\_rx\_ring\_count} neither it nor Spinfer solves anything, while the frontier agent does.

On the development set (Table~\ref{tab:dev}), whose migrations shaped the design, Qwen2.5-Coder solves
\devQwenCoderFull{} of \devBuildable{} with the framework against \devQwenCoderOneShot{} one-shot, Spinfer
\devSpinfer{} and the frontier agent \devFrontier. The general-purpose models again stay near zero.

\begin{table*}[t]
\centering
\caption{Development set, same layout as Table~\ref{tab:test}. These migrations were used to design the framework,
so the numbers are optimistic.}
\label{tab:dev}
\footnotesize
\setlength{\tabcolsep}{4pt}
\begin{tabular}{lrrrrrrrrr}
\toprule
 & & Frontier & & \multicolumn{2}{c}{Qwen2.5-Coder 7B} & \multicolumn{2}{c}{Llama 3.1 8B} & \multicolumn{2}{c}{Gemma 3 4B} \\
\cmidrule(lr){5-6}\cmidrule(lr){7-8}\cmidrule(lr){9-10}
Migration & Bld. & LLM$^\dagger$ & Spinfer & one-shot & \textsc{full} & one-shot & \textsc{full} & one-shot & \textsc{full} \\
\midrule
\texttt{i2c\_probe\_new} & 19 & 19 & 0 & 12.7 [0--19] & 19.0 & 0.0 & 0.0 & 0.0 & 0.0 \\
\texttt{ida\_simple} & 20 & 14 & 9 & 0.0 & 5.3 [5--6] & 0.0 & 1.7 [0--5] & 0.0 & 3.0 [0--9] \\
\texttt{pci\_dma\_wrappers} & 20 & 6 & 2 & 0.0 & 1.3 [0--2] & 0.0 & 0.0 & 0.0 & 0.0 \\
\texttt{platform\_remove\_new} & 20 & 19 & 0 & 0.0 & 0.0 & 0.0 & 0.0 & 0.0 & 0.0 \\
\texttt{strlcpy\_strscpy} & 15 & 14 & 13 & 8.7 [0--13] & 12.3 [11--13] & 0.0 & 0.0 & 0.0 & 0.0 \\
\midrule
\textbf{All} & 94 & 72 & 24 & 21.3 [13--32] & 38.0 [37--39] & 0.0 & 1.7 [0--5] & 0.0 & 3.0 [0--9] \\
\bottomrule
\end{tabular}

\end{table*}

\begin{table*}[t]
\centering
\caption{Ablation: conversions solved as components are added (mean over three seeds, [min--max]). Every
configuration is composed by stage~3; the last row adds the compile critic of stage~4.}
\label{tab:ablation}
\footnotesize
\setlength{\tabcolsep}{4pt}
\begin{tabular}{lrrrrrr}
\toprule
 & \multicolumn{3}{c}{Development (of 94)} & \multicolumn{3}{c}{Test (of 122)} \\
\cmidrule(lr){2-4}\cmidrule(lr){5-7}
Configuration & Qwen2.5-Coder & Llama & Gemma & Qwen2.5-Coder & Llama & Gemma \\
\midrule
\textsc{one-shot} & 21.3 [13--32] & 0.0 & 0.0 & 1.0 & 0.0 & 0.0 \\
+ feedback & 21.3 [13--32] & 0.0 & 0.0 & 1.0 & 2.0 [0--6] & 0.0 \\
+ decomposition & 21.7 [13--32] & 0.0 & 0.0 & 18.0 & 2.0 [0--6] & 0.0 \\
+ guidance & 29.0 [21--34] & 2.3 [0--7] & 0.0 & 35.3 [32--41] & 10.0 [9--12] & 0.0 \\
+ auto-repair & 38.0 [37--39] & 1.7 [0--5] & 3.0 [0--9] & 35.3 [32--41] & 7.0 [0--12] & 2.7 [0--8] \\
+ generality critic & 38.0 [37--39] & 1.7 [0--5] & 3.0 [0--9] & 35.3 [32--41] & 7.0 [0--12] & 2.7 [0--8] \\
+ compile critic (\textsc{full}) & 38.0 [37--39] & 1.7 [0--5] & 3.0 [0--9] & 35.3 [32--41] & 7.0 [0--12] & 2.7 [0--8] \\
\bottomrule
\end{tabular}

\end{table*}

\paragraph{Which components matter.} Table~\ref{tab:ablation} adds one component at a time. For Qwen2.5-Coder on
the test set, test feedback alone changes nothing (\testAblQwenCoderOneShot{} to \testAblQwenCoderFeedback):
asked for one rule covering every variant at once, the model never produces a rule good enough for feedback to
refine. Decomposition into variant leaves is the first large step (\testAblQwenCoderDecomp), and guidance (minimal
edits, context, interfaces, hints, lint and stall handling) roughly doubles it again
(\testAblQwenCoderGuidance). On the development set the order differs: decomposition adds little
(\devAblQwenCoderFeedback{} to \devAblQwenCoderDecomp), guidance and auto-repair add most
(\devAblQwenCoderGuidance, then \devAblQwenCoderRepair). The general-purpose models fail for a more basic reason:
they rarely write SmPL that parses. Without auto-repair, the best rule of \parseFailGemmaNoFix\% of Gemma's leaf
tasks does not parse (\parseFailLlamaNoFix\% for Llama, \parseFailQwenCoderNoFix\% for Qwen2.5-Coder); auto-repair
lowers this to \parseFailGemmaFix\%, \parseFailLlamaFix\% and \parseFailQwenCoderFix\%. For Gemma, auto-repair
is what makes any success possible at all; for Llama on the test set it does not help (\testAblLlamaGuidance{}
before, \testAblLlamaRepair{} after). Adding the generality critic left the mean number of solved conversions
unchanged for every model on both sets. We keep it because it is cheap and did not hurt, but our data give no
evidence that it is needed.

\begin{table*}[t]
\centering
\caption{Precision of the outputs on buildable targets (local models: mean over seeds). \emph{Changed}: targets
the rule edits; \emph{Divergent}: the rule makes an edit the human did not make; \emph{Broken}: the output matches
the human fix (exactly or as a subset) but does not compile; \emph{Precision}: solved / changed.}
\label{tab:safety}
\footnotesize
\setlength{\tabcolsep}{5pt}
\begin{tabular}{llrrrrr}
\toprule
Set & Source & Success & Changed & Divergent & Broken & Precision \\
\midrule
dev & Frontier LLM$^\dagger$ & 72.0 & 91.0 & 17.0 & 0.0 & 0.79 \\
dev & Spinfer & 24.0 & 47.0 & 13.0 & 0.0 & 0.51 \\
dev & Qwen2.5-Coder 7B, before compile critic & 38.0 & 82.7 & 12.0 & 20.0 & 0.46 \\
dev & Qwen2.5-Coder 7B, \textsc{full} & 38.0 & 62.7 & 12.0 & 0.0 & 0.61 \\
dev & Llama 3.1 8B, before compile critic & 1.7 & 22.3 & 10.7 & 0.0 & 0.07 \\
dev & Llama 3.1 8B, \textsc{full} & 1.7 & 20.7 & 8.7 & 0.0 & 0.08 \\
dev & Gemma 3 4B, before compile critic & 3.0 & 9.3 & 2.7 & 0.0 & 0.32 \\
dev & Gemma 3 4B, \textsc{full} & 3.0 & 9.3 & 2.7 & 0.0 & 0.32 \\
\midrule
test & Frontier LLM$^\dagger$ & 66.0 & 94.0 & 28.0 & 0.0 & 0.70 \\
test & Spinfer & 41.0 & 74.0 & 18.0 & 13.0 & 0.55 \\
test & Qwen2.5-Coder 7B, before compile critic & 35.3 & 74.3 & 22.0 & 13.0 & 0.48 \\
test & Qwen2.5-Coder 7B, \textsc{full} & 35.3 & 48.0 & 8.7 & 0.0 & 0.74 \\
test & Llama 3.1 8B, before compile critic & 7.0 & 27.3 & 0.3 & 0.0 & 0.26 \\
test & Llama 3.1 8B, \textsc{full} & 7.0 & 27.3 & 0.3 & 0.0 & 0.26 \\
test & Gemma 3 4B, before compile critic & 2.7 & 34.7 & 0.0 & 12.0 & 0.08 \\
test & Gemma 3 4B, \textsc{full} & 2.7 & 22.7 & 0.0 & 0.0 & 0.12 \\
\bottomrule
\end{tabular}

\end{table*}

\paragraph{The compile critic buys precision at no cost.} The compile critic never changes the number of
solved conversions (Table~\ref{tab:ablation}), but it changes what the rule does to the others
(Table~\ref{tab:safety}). Before it, Qwen2.5-Coder's composed rules on the test set produce
\testSafeQwenCoderPreBroken{} outputs per seed that match the human fix but do not build, and
\testSafeQwenCoderPreDivergent{} divergent ones. The critic either repairs such a rule or abstains for that
migration, leaving \testSafeQwenCoderBroken{} broken and \testSafeQwenCoderDivergent{} divergent outputs and raising
precision from \testSafeQwenCoderPrePrecision{} to \testSafeQwenCoderPrecision{}, above Spinfer
(\testSafeSpinferPrecision, with \testSafeSpinferBroken{} broken outputs) and the frontier agent
(\testSafeFrontierPrecision). On the development set the frontier agent remains more precise
(\devSafeFrontierPrecision{} against \devSafeQwenCoderPrecision). For a maintainer, a rule that converts fewer
files but rarely needs to be undone, or no rule at all, is easier to adopt than one that has to be checked file by
file.

\paragraph{Where the framework fails.} Three failure modes account for the test migrations the local models do not
solve. (i) \emph{Changes that span declarations.} \texttt{gpio\_set\_rv} turns a \texttt{void} callback into one
returning \texttt{int}, returns the value of its last call, and renames the struct field that registers it;
\texttt{ethtool\_rx\_ring\_count} moves a \texttt{switch} case into a new function and registers it in an
\texttt{ethtool\_ops} initialiser. Both need several rules that agree on a function name across a file. The
framework's interface mechanism passes metavariables from an anchor rule to its dependants, but whenever a model
composed a rule for these migrations, its outputs on the training files did not compile (the compiler reports
the registering field renamed while the callback keeps its old type, or a \texttt{return} with a value left in a
\texttt{void} function), repair did not fix them, and the critic abstained
for all three models on every seed. The development
migration \texttt{platform\_remove\_new} has the same shape and the same outcome. (ii) \emph{Over-specific
patterns.} On \texttt{iio\_push\_with\_ts} Qwen2.5-Coder's rule requires the timestamp argument to be a function
call, so it misses the common \texttt{pf->timestamp} form, and it writes \texttt{sizeof(ID2)} where
\texttt{ID2} is the expression \texttt{\&scan}: the output computes the size of a pointer. That output compiles,
which no compile check can catch; it shows up only as a divergent edit. (iii) \emph{Invalid SmPL} remains the main
failure of the two general-purpose models, as quantified above.

\begin{table}[t]
\centering
\caption{Cost per migration of the \textsc{full} configuration on one laptop GPU (RTX~5060, 8\,GB): model calls
and seconds of generation in stage~2, and repair calls of the compile critic (stage~4), mean over migrations and
seeds. Compilation time is excluded.}
\label{tab:cost}
\resizebox{\columnwidth}{!}{\begin{tabular}{lrrrrrr}
\toprule
 & \multicolumn{3}{c}{Development} & \multicolumn{3}{c}{Test} \\
\cmidrule(lr){2-4}\cmidrule(lr){5-7}
Model & Calls & Seconds & Repairs & Calls & Seconds & Repairs \\
\midrule
Qwen2.5-Coder 7B & 13.2 & 21 & 0.4 & 9.7 & 15 & 0.5 \\
Llama 3.1 8B & 16.3 & 31 & 0.1 & 9.7 & 23 & 0.0 \\
Gemma 3 4B & 17.4 & 30 & 0.0 & 9.7 & 16 & 0.3 \\
\bottomrule
\end{tabular}
}
\end{table}

\paragraph{Cost.} Table~\ref{tab:cost} reports generation cost. A migration needs ten to twenty model calls and well
under a minute of generation (\testCostQwenCoderCalls{} calls and \testCostQwenCoderSeconds{}\,s for
Qwen2.5-Coder on the test set). Stage~4 additionally configures a kernel tree at each training commit it checks
and compiles single object files there; build results are cached by commit, file and content.
